% TEMPLATE EXEMPLAR of a reading-guide chapter: a short narrative that leads into paper cards. Cards say what a
% paper does (Idea), what it showed (Results, with the authors' numbers attributed to them), why it matters and
% its limits. Delete this chapter when the book's own literature chapters exist.
%
% arXiv papers use papercard; register them first (key, arXiv id, tier A/B/C, section code, short name):
%   printf 'vaswani2017attention\t1706.03762\tA\tfound\tTransformer\n' > /tmp/rows.tsv
%   python3 -I tools/add_arxiv.py /tmp/rows.tsv
% and then
%   \begin{papercard}{vaswani2017attention}{Transformer}
%   \cl{Idea.} ...  \cl{Results.} ...  \cl{Why it matters.} ...  \cl{Limits.} ...
%   \end{papercard}
% Sources without arXiv metadata use webcard and an entry in refs_manual.bib (checked before it is added):
%   \begin{webcard}{bibkey}{Display name}{Full title}{meta line}{tier}
\chapter{Reading the Literature}
\label{ch:reading}

A reading guide is only as useful as its advice on how to read. This chapter shows the two kinds of card the book
uses and the method behind them: each card tells you what a work does, why it matters and what to read it for, so
that you can decide in a minute whether to open it. The tier, shown by the stars and by the shape of the box, says
how much of the field rests on it.

\section{How to read a paper}
\label{sec:reading-how}

Most papers do not need to be read from the first line to the last. A short article by Keshav turned that
observation into a method.

\begin{webcard}{keshav2007read}{Three-pass reading}{How to Read a Paper}{S.~Keshav \textperiodcentered\ 2007 (ACM SIGCOMM Computer Communication Review)}{A}
\cl{Idea.} Read a paper in up to three passes. The first, of five to ten minutes, gives a bird's-eye view: the title,
abstract and introduction, the section headings and the conclusions. The second, of up to an hour, takes in the
content but skips details such as proofs. The third tries to ``virtually re-implement the paper'', making the
authors' assumptions and re-creating the work.
\cl{Why it matters.} It turns ``read everything'' into a decision after each pass, and the first pass ends with five
questions (category, context, correctness, contributions, clarity) that make a good skeleton for a reading note.
\cl{Limits.} The time budgets are the author's experience, not measurements: he puts the third pass at four or five
hours for a beginner and about an hour for an experienced reader.
\end{webcard}

\section{Where this book's look comes from}
\label{sec:reading-look}

\begin{webcard}{chen2024napkin}{The Napkin}{An Infinitely Large Napkin}{Evan Chen \textperiodcentered\ project page, version 1.6}{C}
\cl{Idea.} In its author's words, ``a light and mostly self-contained introduction to higher math'', from the
undergraduate syllabus to first-year graduate topics, that gives a bird's-eye view of the main ideas of each field.
\cl{Why it matters here.} Its typography (bold sans-serif headings, crimson numerals, a blue box for what matters
most, a rust box for examples, a green bar for remarks) is the model for the look of this book, which uses none of
its text or artwork.
\end{webcard}

\begin{readingbox}
Read the first pass of every Essential card's paper before the second pass of any of them: the map of the field
comes first.
\end{readingbox}
